% -*- mode:LaTex; mode:visual-line; mode:flyspell; fill-column:75-*-

% Special indentation for abstract.
\setlength{\parskip}{1em}
\setlength{\parindent}{0em}

\noindent
Bio-inspired aquatic robots offer a promising route to agile and efficient locomotion in fluid environments, where conventional rigid systems remain limited. In bio-inspred aquatic systems, locomotion is not determined by actuation or control alone, but instead emerges from tightly coupled interactions among body morphology, distributed compliance, actuation, onboard sensing, and the surrounding flow, making analysis, design, and optimization fundamentally challenging. This thesis studies these coupled problems through the development of embodied aquatic robot platforms, variational dynamical models, and integrated sensing strategies, with particular emphasis on undulatory and flapping swimming, variable-stiffness design, strongly coupled fluid–robot interaction, and flow sensing and feedback based on local pressure measurements. The goal is to understand how morphology, dynamics, and sensing can be jointly exploited within a continuous-time dynamical framework to realize robotic swimmers that are more efficient, maneuverable, and adaptive, while also enabling reliable optimization of physical design, sensing configuration, and control through optimization-based fluid dynamic formulations and high-fidelity gradients. These ideas are further examined and validated through experiments on physical robotic platforms.